\chapter{C++ 语法与 STL 速查}\label{ch:04}

\section{本章使用方式}\label{sec:04-1}

这一章放常用的 C++ 写法。忘了容器怎么声明、循环怎么写，或者拿不准整数除法和取模的结果时，可以从这里查：

\begin{itemize}
  \item 一个程序应该怎么开头。
  \item 输入输出怎么写。
  \item 什么时候必须用 \Code{long long}。
  \item 整数除法、取模、向上取整、向下取整怎么写。
  \item 数组、\Code{vector}、\Code{string}、结构体怎么用。
  \item \Code{sort}、\Code{map}、\Code{set}、\Code{queue}、\Code{priority\_queue} 怎么写。
  \item lambda 比较器怎么写，常见错误是什么。
\end{itemize}

写法熟悉就继续用，不必为了把代码写短而临时换一种。

\section{语言标准与代码拼装}\label{sec:04-2}
本书代码使用 \textbf{GNU C++17}。结构化绑定 \Code{auto [x,y]} 和 \Code{std::gcd} 需要 C++17，编译时可使用 \Code{-std=c++17}。

若只支持旧标准，把 \Code{auto [x,y]=p;} 改成 \Code{auto x=p.first; auto y=p.second;}；\Code{gcd} 可自行写欧几里得算法。\Code{bits/stdc++.h} 是 GNU 扩展；非 GNU 环境改用对应标准头文件。

\textbf{模板不是一整份提交代码。}先写下面的主函数骨架；结构体、普通函数放在 \Code{main} 外，输入、循环、输出放在 \Code{main} 或 \Code{solve} 内。不要把两个同名 \Code{main}、\Code{Node}、\Code{Edge} 或 \Code{INF} 一起粘贴。带 \Code{check}、\Code{handle...} 和省略号的代码是待按题意补全的框架。

\section{最小程序骨架}\label{sec:04-3}

大部分 CSP 题目都可以从这个骨架开始：

\begin{lstlisting}
#include <bits/stdc++.h>
using namespace std;

using ll = long long;

int main() {
    ios::sync_with_stdio(false);
    cin.tie(nullptr);

    // write code here

    return 0;
}
\end{lstlisting}

说明：

\begin{itemize}
  \item \Code{\#include <bits/stdc++.h>} 是 GNU C++ 的扩展头文件，包含大多数常用标准库。
  \item \Code{using namespace std;} 可以少写 \Code{std::}。
  \item \Code{using ll = long long;} 以后可以用 \Code{ll} 表示 \Code{long long}。
  \item \Code{ios::sync\_with\_stdio(false); cin.tie(nullptr);} 用来加速 C++ 输入输出。
  \item 使用上面两行加速后，不要混用 \Code{scanf/printf} 和 \Code{cin/cout}。
\end{itemize}

\section{输入输出}\label{sec:04-4}

\subsection{普通输入}

\begin{lstlisting}
int n;
cin >> n;

int a, b;
cin >> a >> b;

string s;
cin >> s; // reads until whitespace
\end{lstlisting}

\subsection{读入数组}

\begin{lstlisting}
int n;
cin >> n;
vector<int> a(n);
for (int i = 0; i < n; i++) {
    cin >> a[i];
}
\end{lstlisting}

如果题目下标从 1 开始，考场上可以直接开成 \Code{n + 1}：

\begin{lstlisting}
int n;
cin >> n;
vector<int> a(n + 1);
for (int i = 1; i <= n; i++) {
    cin >> a[i];
}
\end{lstlisting}

\subsection{多组数据}

题目明确给出 $T$：

\begin{lstlisting}
int T;
cin >> T;
while (T--) {
    solve();
}
\end{lstlisting}

题目没有给 $T$，读到文件结束：

\begin{lstlisting}
int a, b;
while (cin >> a >> b) {
    cout << a + b << '\n';
}
\end{lstlisting}

\Warn{多组数据常见坑：}每组开始前要清空数组、图、队列、\Code{map}、答案变量。

\subsection{整行输入}

\Code{cin >> s} 只能读到空格前。如果要读整行，用 \Code{getline}：

\begin{lstlisting}
string line;
getline(cin, line);
\end{lstlisting}

如果前面刚用了 \Code{cin >> n}，再用 \Code{getline}，要先吃掉换行：

\begin{lstlisting}
int n;
cin >> n;
cin.ignore(numeric_limits<streamsize>::max(), '\n');

string line;
getline(cin, line);
\end{lstlisting}

如果担心行尾有多个字符，可以写：

\begin{lstlisting}
cin.ignore(numeric_limits<streamsize>::max(), '\n');
\end{lstlisting}

\subsection{输出}

\begin{lstlisting}
cout << ans << '\n';
cout << a << ' ' << b << '\n';
\end{lstlisting}

\Warn{不要输出提示语。}例如不要写 \Code{cout << "answer=" << ans;}，除非题目明确要求。

\subsection{小数输出}

\begin{lstlisting}
double x = 3.1415926;
cout << fixed << setprecision(2) << x << '\n'; // 3.14
\end{lstlisting}

说明：

\begin{itemize}
  \item \Code{fixed} 表示按小数位输出。
  \item \Code{setprecision(2)} 表示保留 2 位小数。
  \item 需要包含 \Code{<iomanip>}，但 \Code{bits/stdc++.h} 已包含。
\end{itemize}

\section{数据类型}\label{sec:04-5}

\subsection{整数类型}

\begin{longtable}{p{3cm}p{4cm}p{7cm}}
\toprule
类型 & 大约范围 & 使用场景 \\
\midrule
\endfirsthead
\toprule
类型 & 大约范围 & 使用场景 \\
\midrule
\endhead
\Code{int} & $-2.1 \times 10^9$ 到 $2.1 \times 10^9$ & 下标、普通计数、小范围整数。 \\
\Code{long long} & $-9.2 \times 10^{18}$ 到 $9.2 \times 10^{18}$ & 求和、乘法、答案、距离、方案数。 \\
\Code{unsigned long long} & 0 到约 $1.8 \times 10^{19}$ & 很少必须用，新手慎用。 \\
\bottomrule
\end{longtable}

必须用 \Code{long long} 的常见情况：

\begin{itemize}
  \item $n$ 个数求和，每个数可达 $10^9$。
  \item 两个最大可达 $10^9$ 的数相乘。
  \item 最短路边权和可能很大。
  \item 方案数、组合数、累计时间、累计代价。
\end{itemize}

乘法防溢出：

\begin{lstlisting}
int a = 1000000000;
int b = 1000000000;
long long x = 1LL * a * b;
\end{lstlisting}

\Warn{错误写法：}

\begin{lstlisting}
long long x = a * b; // a*b already overflowed as int
\end{lstlisting}

\subsection{浮点类型}

\begin{itemize}
  \item \Code{double} 是常用浮点数。
  \item 判断浮点相等时不要直接用 \Code{==}，通常用误差。
\end{itemize}

\begin{lstlisting}
const double EPS = 1e-9;
if (abs(a - b) < EPS) {
    // a and b are almost equal
}
\end{lstlisting}

\subsection{字符和布尔}

\begin{lstlisting}
char c = 'A';
bool ok = true;
\end{lstlisting}

字符判断：

\begin{lstlisting}
isdigit((unsigned char)c); // is digit
isalpha((unsigned char)c); // is letter
islower((unsigned char)c);
isupper((unsigned char)c);
tolower((unsigned char)c);
toupper((unsigned char)c);
\end{lstlisting}

\Warn{注意：}\Code{isdigit((unsigned char)c)} 返回非零表示真，不一定返回 1。

\section{表达式、整除和取整}\label{sec:04-6}

\subsection{基本运算符}

C++ 中常用算术运算符如下：

\begin{longtable}{p{3cm}p{4cm}p{7cm}}
\toprule
运算符 & 含义 & 注意点 \\
\midrule
\endfirsthead
\toprule
运算符 & 含义 & 注意点 \\
\midrule
\endhead
\Code{+ - *} & 加、减、乘 & 两个 \Code{int} 相乘仍先按 \Code{int} 计算。 \\
\Code{/} & 除法 & 两边都是整数时做整数除法，结果向 0 截断。 \\
\Code{\%} & 取余 & 只能用于整数；负数取余的结果可能为负。 \\
\Code{+= -= *= /= \%=} & 复合赋值 & \Code{x += y} 等价于 \Code{x = x + y}。 \\
\Code{++ --} & 自增、自减 & 简单循环中优先写 \Code{i++} 或 \Code{++i}，不要嵌进复杂表达式。 \\
\bottomrule
\end{longtable}

\Warn{整数除法不是数学里的向下取整。}在 C++ 中，\Code{-3 / 2} 的结果是 \Code{-1}，因为整数除法向 0 截断；数学意义上的 $\lfloor -3/2 \rfloor$ 是 $-2$。

\subsection{整数除法和类型转换}

如果两边都是整数，\Code{/} 会丢掉小数部分：

\begin{lstlisting}
int a = 5, b = 2;
cout << a / b << '\n'; // 2
\end{lstlisting}

如果要小数结果，至少把一个操作数转成浮点数：

\begin{lstlisting}
cout << 1.0 * a / b << '\n';       // 2.5
cout << (double)a / b << '\n';     // 2.5
\end{lstlisting}

整数答案尽量用整数运算求。大整数转成 \Code{double} 后可能丢失低位，取整时尤其要留意。

\subsection{正整数向上取整}

最常见场景：$a$ 个物品，每组最多放 $b$ 个，问需要几组。若 $a \ge 0, b > 0$：

\begin{lstlisting}
long long ceil_div_positive(long long a, long long b) {
    return a / b + (a % b != 0);
}
\end{lstlisting}

常见特例：

\begin{lstlisting}
long long half_up = (n + 1) / 2;          // ceil(n / 2), n >= 0
long long blocks = (n + block - 1) / block;
\end{lstlisting}

\Warn{这个简式只适合 $a \ge 0, b > 0$。}如果 $a$ 可能为负，或者 $b$ 可能为负，不能直接套 \Code{a/b+(a\%b!=0)}。

\subsection{有符号整数向下和向上取整}

如果题目涉及负坐标、负偏移、时间差、数学分段，建议使用下面的通用函数。它们要求 \Code{b != 0}，并且遵循数学意义上的 floor 和 ceil。

\begin{lstlisting}
// Require b != 0; exclude LLONG_MIN / -1.
long long floor_div(long long a, long long b) {
    long long q = a / b;
    long long r = a % b;
    if (r != 0 && ((r > 0) != (b > 0))) q--;
    return q;
}

long long ceil_div(long long a, long long b) {
    long long q = a / b;
    long long r = a % b;
    if (r != 0 && ((r > 0) == (b > 0))) q++;
    return q;
}
\end{lstlisting}

快速自测：

\begin{longtable}{p{4cm}p{4cm}p{4cm}}
\toprule
表达式 & \Code{floor\_div} & \Code{ceil\_div} \\
\midrule
\endfirsthead
\toprule
表达式 & \Code{floor\_div} & \Code{ceil\_div} \\
\midrule
\endhead
$3 / 2$ & 1 & 2 \\
$-3 / 2$ & -2 & -1 \\
$3 / -2$ & -2 & -1 \\
$-3 / -2$ & 1 & 2 \\
\bottomrule
\end{longtable}

\subsection{负数取模}

C++ 中 \Code{x \% m} 的符号跟 \Code{x} 相同。例如 \Code{-1 \% 5} 是 \Code{-1}，不是 \Code{4}。如果题目需要 $0$ 到 $m-1$ 的标准余数，写：

\begin{lstlisting}
long long norm_mod(long long x, long long m) {
    long long r = x % m;
    if (r < 0) r += m;
    return r;
}
\end{lstlisting}

适用：\Code{m > 0}。若 \Code{x} 可能小于 \Code{-m} 很多，上面仍然正确，因为 C++ 的一次取余后 \Code{r} 已经在 \Code{(-m,m)} 内。

\subsection{运算优先级和括号}

常用的优先级关系如下。混合运算拿不准时，可以多加一层括号：

\begin{itemize}
  \item 乘、除、取模优先于加减。
  \item 比较运算优先级低于加减乘除。
  \item \Code{\&\&} 优先于 \Code{||}。
  \item 位运算和比较混在一起时必须加括号。
\end{itemize}

\begin{lstlisting}
if ((mask & (1 << k)) != 0) {
    // bit k is 1
}
\end{lstlisting}

\Warn{有疑问就加括号。}括号不会让程序变慢，却能显著减少读错表达式的风险。

\subsection{常量和初始化}

常用写法：

\begin{lstlisting}
const int N = 100000 + 5;
const long long INF = (long long)4e18;
const int MOD = 1000000007;

int cnt = 0;
long long sum = 0;
bool ok = true;
\end{lstlisting}

提醒：

\begin{itemize}
  \item 局部普通变量不会自动清零，定义后要赋初值。
  \item 全局数组会默认清零，但多组数据不会自动恢复初始状态。
  \item \Code{const} 常量不能被修改，适合写数组上限、模数、无穷大。
\end{itemize}


\subsection{四舍五入与银行家舍入不能混用}
\Code{llround} 的半整数规则是远离 0；题面要求“恰好一半时取偶数”时应自行实现。若输入为非负数、恰有一位小数，且整数部分不超过 $10^5$，可以完全避免浮点：
\begin{lstlisting}
pair<long long,long long> roundOneDecimal(const string& s) {
    size_t dot = s.find('.');
    long long a = stoll(s.substr(0, dot));
    int b = s[dot + 1] - '0';
    long long halfUp = a + (b >= 5);
    long long halfEven = a + (b > 5 || (b == 5 && a % 2 == 1));
    return {halfUp, halfEven};
}
\end{lstlisting}
例如 2.5 分别得到 3 与 2，3.5 两者均为 4。此代码仅针对上述输入格式；不能直接用于负数、多位小数或超大整数。\Code{setprecision} 负责显示格式，不替代题目定义的舍入算法。

\section{控制流}\label{sec:04-7}

\subsection{if}

\begin{lstlisting}
if (x > 0) {
    cout << "positive\n";
} else if (x == 0) {
    cout << "zero\n";
} else {
    cout << "negative\n";
}
\end{lstlisting}

\subsection{for}

\begin{lstlisting}
for (int i = 0; i < n; i++) {
    // 0, 1, ..., n-1
}

for (int i = 1; i <= n; i++) {
    // 1, 2, ..., n
}
\end{lstlisting}

倒序：

\begin{lstlisting}
for (int i = n - 1; i >= 0; i--) {
    // be careful when i is unsigned
}
\end{lstlisting}

\subsection{while}

\begin{lstlisting}
while (condition) {
    // loop
}
\end{lstlisting}

读到结束：

\begin{lstlisting}
int x;
while (cin >> x) {
    // process x
}
\end{lstlisting}

\subsection{break 和 continue}

\begin{lstlisting}
for (int i = 0; i < n; i++) {
    if (a[i] < 0) continue; // skip this round
    if (a[i] == target) break; // exit loop
}
\end{lstlisting}

\section{auto、引用和范围 for}\label{sec:04-8}

\subsection{auto}

\Code{auto} 让编译器根据右侧表达式推断类型，常用于迭代器、pair、复杂容器遍历。

\begin{lstlisting}
auto it = lower_bound(a.begin(), a.end(), x);
auto p = make_pair(3, 5);
\end{lstlisting}

遍历 \Code{map}：

\begin{lstlisting}
map<string, int> cnt;
for (auto [key, value] : cnt) {
    cout << key << ' ' << value << '\n';
}
\end{lstlisting}

\Warn{注意复制。}\Code{auto [key, value]} 会复制键和值；如果元素很大，写 \Code{const auto\&} 更合适。

\begin{lstlisting}
for (const auto& item : records) {
    // read item without copying
}
\end{lstlisting}

\subsection{引用}

引用相当于变量的别名。函数传大对象时常用引用避免复制。

\begin{lstlisting}
void printVector(const vector<int>& a) {
    for (int x : a) cout << x << ' ';
}

void addOne(vector<int>& a) {
    for (int& x : a) x++;
}
\end{lstlisting}

区别：

\begin{itemize}
  \item \Code{const vector<int>\& a}：不复制，也不允许函数内修改。
  \item \Code{vector<int>\& a}：不复制，允许函数内修改原数组。
  \item \Code{vector<int> a}：会复制整个数组，数据大时慢。
\end{itemize}

\subsection{范围 for}

只读遍历：

\begin{lstlisting}
for (int x : a) {
    cout << x << '\n';
}
\end{lstlisting}

修改原元素：

\begin{lstlisting}
for (int& x : a) {
    x++;
}
\end{lstlisting}

\Warn{如果忘了写引用，修改不会回到原数组。}

\begin{lstlisting}
for (int x : a) {
    x++; // only modifies the copy
}
\end{lstlisting}

\subsection{size 的类型}

\Code{a.size()} 的返回类型是无符号整数。考场上为了减少类型麻烦，常写成：

\begin{lstlisting}
int n = (int)a.size();
for (int i = 0; i < n; i++) {
    // use a[i]
}
\end{lstlisting}

倒序遍历不要写下面这种：

\begin{lstlisting}
for (int i = a.size() - 1; i >= 0; i--) {
    // dangerous when a is empty
}
\end{lstlisting}

更稳的写法：

\begin{lstlisting}
for (int i = (int)a.size() - 1; i >= 0; i--) {
    // ok even when a is empty: i starts at -1
}
\end{lstlisting}

\section{函数}\label{sec:04-9}

把重复逻辑写成函数，能减少大模拟和算法题的错误。

\begin{lstlisting}
int add(int a, int b) {
    return a + b;
}

bool isEven(int x) {
    return x % 2 == 0;
}
\end{lstlisting}

传 vector：

\begin{lstlisting}
long long sumVector(const vector<int>& a) {
    long long sum = 0;
    for (int x : a) sum += x;
    return sum;
}
\end{lstlisting}

说明：

\begin{itemize}
  \item \Code{const vector<int>\& a} 表示不复制整个数组，并且函数内部不修改它。
  \item 如果函数需要修改数组，去掉 \Code{const}。
\end{itemize}

\begin{lstlisting}
void addOne(vector<int>& a) {
    for (int& x : a) x++;
}
\end{lstlisting}

\section{数组}\label{sec:04-10}

\subsection{普通数组}

全局数组适合大数组：

\begin{lstlisting}
const int N = 100000 + 5;
int a[N];
long long pre[N];
\end{lstlisting}

说明：

\begin{itemize}
  \item 全局数组默认初始化为 0。
  \item 局部大数组可能爆栈，尽量放全局或用 \Code{vector}。
\end{itemize}

\subsection{二维数组}

\begin{lstlisting}
const int N = 1005;
int grid[N][N];
\end{lstlisting}

如果规模不确定：

\begin{lstlisting}
int n, m;
cin >> n >> m;
vector<vector<int>> grid(n, vector<int>(m, 0));
\end{lstlisting}

\Warn{内存估算：}\Code{int a[10000][10000]} 大约 400MB，通常危险。

\subsection{初始化}

全局普通数组默认是 0。如果要重置数组：

\begin{lstlisting}
int a[1005];
memset(a, 0, sizeof a);
memset(a, -1, sizeof a);
\end{lstlisting}

\Warn{注意：}\Code{memset} 适合把数组设为 0 或 -1，不适合把 \Code{int} 数组设成 1、2、无穷大。

设置成大数：

\begin{lstlisting}
const int INF = 0x3f3f3f3f;
int dist[1005];
memset(dist, 0x3f, sizeof dist);
\end{lstlisting}

\Code{vector} 初始化：

\begin{lstlisting}
vector<int> a(n, 0);
vector<long long> dist(n + 1, (long long)4e18);
vector<vector<int>> grid(n, vector<int>(m, 0));
\end{lstlisting}

\section{vector}\label{sec:04-11}

\subsection{创建}

\begin{lstlisting}
vector<int> a;          // empty
vector<int> b(10);      // 10 zeros
vector<int> c(10, -1);  // 10 values, all -1
\end{lstlisting}

\subsection{常用操作}

\begin{lstlisting}
a.push_back(5);
a.pop_back();
a.clear();
int n = (int)a.size();
bool empty = a.empty();
\end{lstlisting}

访问：

\begin{lstlisting}
cout << a[0] << '\n';
cout << a.back() << '\n';
\end{lstlisting}

\Warn{访问前检查非空。}\Code{a.back()}、\Code{a[0]} 在空数组上会出错。

\subsection{遍历}

\begin{lstlisting}
for (int i = 0; i < (int)a.size(); i++) {
    cout << a[i] << '\n';
}

for (int x : a) {
    cout << x << '\n';
}

for (int& x : a) {
    x++; // modify original element
}
\end{lstlisting}

\subsection{删除元素}

删除最后一个：

\begin{lstlisting}
a.pop_back();
\end{lstlisting}

删除某个位置：

\begin{lstlisting}
a.erase(a.begin() + pos);
\end{lstlisting}

\Warn{注意：}\Code{erase} 是 $O(n)$，大量删除时可能超时。

删除所有等于 x 的元素：

\begin{lstlisting}
a.erase(remove(a.begin(), a.end(), x), a.end());
\end{lstlisting}

\subsection{二维 vector}

二维网格：

\begin{lstlisting}
int n, m;
cin >> n >> m;
vector<vector<int>> a(n, vector<int>(m));
for (int i = 0; i < n; i++) {
    for (int j = 0; j < m; j++) {
        cin >> a[i][j];
    }
}
\end{lstlisting}

二维字符图：

\begin{lstlisting}
int n, m;
cin >> n >> m;
vector<string> g(n);
for (int i = 0; i < n; i++) {
    cin >> g[i];
}
cout << g[0][0] << '\n';
\end{lstlisting}

图的邻接表：

\begin{lstlisting}
int n, m;
cin >> n >> m;
vector<vector<int>> g(n + 1);
for (int i = 0; i < m; i++) {
    int u, v;
    cin >> u >> v;
    g[u].push_back(v);
    g[v].push_back(u);
}
\end{lstlisting}

\section{string}\label{sec:04-12}

\subsection{基本操作}

\begin{lstlisting}
string s = "abc";
int n = (int)s.size();
char c = s[0];
s.push_back('d');
s += "ef";
\end{lstlisting}

\subsection{子串}

\begin{lstlisting}
string s = "abcdef";
string t = s.substr(2, 3); // "cde"
\end{lstlisting}

\Code{substr(pos, len)} 从 \Code{pos} 开始取 \Code{len} 个字符，下标从 0 开始。

\subsection{查找}

\begin{lstlisting}
string s = "abcdef";
if (s.find("cd") != string::npos) {
    cout << "found\n";
}
\end{lstlisting}

\Warn{不要写：}

\begin{lstlisting}
if (s.find("cd")) {
    // wrong when position is 0
}
\end{lstlisting}

\subsection{字符串转数字}

\begin{lstlisting}
int x = stoi("123");
long long y = stoll("1234567890123");
double z = stod("3.14");
\end{lstlisting}

数字转字符串：

\begin{lstlisting}
string s = to_string(123);
\end{lstlisting}

\Warn{如果数字很长，}\Code{stoi/stoll} 会溢出或异常，要用字符串模拟。

\subsection{按空格分割}

\begin{lstlisting}
string line;
getline(cin, line);
stringstream ss(line);

string word;
while (ss >> word) {
    cout << word << '\n';
}
\end{lstlisting}

\section{pair 和 tuple}\label{sec:04-13}

\subsection{pair}

\begin{lstlisting}
pair<int, int> p = {3, 5};
cout << p.first << ' ' << p.second << '\n';
\end{lstlisting}

常用于坐标、边、距离和编号：

\begin{lstlisting}
vector<pair<int, int>> points;
points.push_back({x, y});
\end{lstlisting}

默认排序：

\begin{lstlisting}
sort(points.begin(), points.end());
\end{lstlisting}

默认先按 \Code{first} 升序，再按 \Code{second} 升序。

\subsection{tuple}

字段超过两个时可以用 \Code{tuple}，但新手更推荐结构体。

\begin{lstlisting}
tuple<int, int, int> t = {1, 2, 3};
auto [a, b, c] = t;
\end{lstlisting}

\section{结构体 struct}\label{sec:04-14}

结构体适合表示学生、任务、边、事件、状态。

\begin{lstlisting}
struct Student {
    int id;
    string name;
    int score;
};

Student s;
s.id = 1;
s.name = "Alice";
s.score = 90;
\end{lstlisting}

数组：

\begin{lstlisting}
vector<Student> students;
students.push_back({1, "Alice", 90});
\end{lstlisting}

图的边：

\begin{lstlisting}
struct Edge {
    int to;
    int w;
};

vector<vector<Edge>> g(n + 1);
g[u].push_back({v, w});
\end{lstlisting}

\section{排序 sort}\label{sec:04-15}

\subsection{普通排序}

\begin{lstlisting}
sort(a.begin(), a.end()); // ascending
sort(a.rbegin(), a.rend()); // descending
\end{lstlisting}

数组排序：

\begin{lstlisting}
sort(a, a + n);
\end{lstlisting}

\subsection{自定义排序：从大到小}

\begin{lstlisting}
sort(a.begin(), a.end(), [](int x, int y) {
    return x > y;
});
\end{lstlisting}

\subsection{结构体排序}

成绩高的在前；成绩相同，编号小的在前：

\begin{lstlisting}
struct Student {
    int id;
    int score;
};

vector<Student> v;

sort(v.begin(), v.end(), [](const Student& a, const Student& b) {
    if (a.score != b.score) return a.score > b.score;
    return a.id < b.id;
});
\end{lstlisting}

\subsection{按 vector 的某一列排序}

\begin{lstlisting}
vector<vector<int>> a;

sort(a.begin(), a.end(), [](const vector<int>& x, const vector<int>& y) {
    return x[0] < y[0]; // sort by column 0 ascending
});
\end{lstlisting}

先按第 0 列升序，再按第 1 列降序：

\begin{lstlisting}
sort(a.begin(), a.end(), [](const vector<int>& x, const vector<int>& y) {
    if (x[0] != y[0]) return x[0] < y[0];
    return x[1] > y[1];
});
\end{lstlisting}

\subsection{stable\_sort}

如果题目要求相同关键字保持原输入顺序：

\begin{lstlisting}
stable_sort(v.begin(), v.end(), [](const Student& a, const Student& b) {
    return a.score > b.score;
});
\end{lstlisting}

\subsection{比较器高危错误}

比较器必须表达“谁应该排在前面”。

\Good{正确：}

\begin{lstlisting}
sort(v.begin(), v.end(), [](const Student& a, const Student& b) {
    if (a.score != b.score) return a.score > b.score;
    return a.id < b.id;
});
\end{lstlisting}

\Warn{错误：}

\begin{lstlisting}
sort(v.begin(), v.end(), [](const Student& a, const Student& b) {
    return a.score >= b.score; // wrong: equality should return false
});
\end{lstlisting}

相等时返回 \Code{true} 可能导致排序行为错误。

\section{二分相关函数}\label{sec:04-16}

\subsection{binary\_search}

\begin{lstlisting}
sort(a.begin(), a.end());
bool ok = binary_search(a.begin(), a.end(), x);
\end{lstlisting}

\subsection{lower\_bound 和 upper\_bound}

\begin{lstlisting}
sort(a.begin(), a.end());

auto it1 = lower_bound(a.begin(), a.end(), x); // first >= x
auto it2 = upper_bound(a.begin(), a.end(), x); // first > x
\end{lstlisting}

下标：

\begin{lstlisting}
int pos = lower_bound(a.begin(), a.end(), x) - a.begin();
\end{lstlisting}

\Warn{解引用前检查：}

\begin{lstlisting}
auto it = lower_bound(a.begin(), a.end(), x);
if (it != a.end()) {
    cout << *it << '\n';
}
\end{lstlisting}

\section{map 和 unordered\_map}\label{sec:04-17}

\subsection{map}

\Code{map} 按 key 自动升序排列，常用于字符串计数、编号映射、需要有序遍历的场景。

\begin{lstlisting}
map<string, int> cnt;
cnt["apple"]++;
cnt["banana"] += 2;

cout << cnt["apple"] << '\n';
\end{lstlisting}

判断是否存在：

\begin{lstlisting}
if (cnt.count("apple")) {
    cout << "exists\n";
}

auto it = cnt.find("apple");
if (it != cnt.end()) {
    cout << it->second << '\n';
}
\end{lstlisting}

\Warn{注意：}\Code{cnt[key]} 会在 key 不存在时自动创建一个默认值。

遍历：

\begin{lstlisting}
for (auto [key, value] : cnt) {
    cout << key << ' ' << value << '\n';
}
\end{lstlisting}

\subsection{unordered\_map}

\Code{unordered\_map} 平均更快，但没有顺序。

\begin{lstlisting}
unordered_map<string, int> cnt;
cnt["apple"]++;
\end{lstlisting}

考场建议：

\begin{itemize}
  \item 需要按 key 排序输出，用 \Code{map}。
  \item 只做计数且数据量大，用 \Code{unordered\_map}。
  \item 自定义 key 时新手优先考虑转成字符串或 \Code{long long} 编码，不要现场写复杂 hash。
\end{itemize}

\section{unordered\_set}\label{sec:04-18}

\Code{unordered\_set} 用来快速判断一个值是否出现过，不关心顺序。

\begin{lstlisting}
unordered_set<int> seen;
seen.insert(5);

if (seen.count(5)) {
    cout << "yes\n";
}
\end{lstlisting}

字符串集合：

\begin{lstlisting}
unordered_set<string> words;
words.insert("abc");
if (words.count("abc")) {
    cout << "exists\n";
}
\end{lstlisting}

如果需要按升序遍历，用 \Code{set}，不要用 \Code{unordered\_set}。

\section{set 和 multiset}\label{sec:04-19}

\subsection{set}

\Code{set} 自动排序且不重复。

\begin{lstlisting}
set<int> s;
s.insert(3);
s.insert(1);
s.erase(3);

if (s.count(1)) {
    cout << "exists\n";
}
\end{lstlisting}

查找第一个大于等于 x 的元素：

\begin{lstlisting}
auto it = s.lower_bound(x);
if (it != s.end()) {
    cout << *it << '\n';
}
\end{lstlisting}

\subsection{multiset}

\Code{multiset} 允许重复元素。

\begin{lstlisting}
multiset<int> s;
s.insert(5);
s.insert(5);
cout << s.count(5) << '\n'; // 2
\end{lstlisting}

删除一个 5：

\begin{lstlisting}
auto it = s.find(5);
if (it != s.end()) s.erase(it);
\end{lstlisting}

删除所有 5：

\begin{lstlisting}
s.erase(5);
\end{lstlisting}

\section{queue、stack、deque}\label{sec:04-20}

\subsection{queue}

BFS 常用队列。

\begin{lstlisting}
queue<int> q;
q.push(1);
int x = q.front();
q.pop();
bool empty = q.empty();
\end{lstlisting}

\subsection{stack}

括号匹配、表达式、DFS 模拟可用栈。

\begin{lstlisting}
stack<char> st;
st.push('(');
char c = st.top();
st.pop();
\end{lstlisting}

\Warn{调用 \Code{front/top} 前必须检查非空。}

\subsection{deque}

双端队列，两边都能进出。

\begin{lstlisting}
deque<int> dq;
dq.push_back(1);
dq.push_front(2);
dq.pop_back();
dq.pop_front();
\end{lstlisting}

单调队列会用到 \Code{deque}。

\section{priority\_queue}\label{sec:04-21}

\subsection{大根堆}

默认是大根堆，最大值在顶部。

\begin{lstlisting}
priority_queue<int> pq;
pq.push(3);
pq.push(10);
cout << pq.top() << '\n'; // 10
pq.pop();
\end{lstlisting}

\subsection{小根堆}

\begin{lstlisting}
priority_queue<int, vector<int>, greater<int>> pq;
pq.push(3);
pq.push(10);
cout << pq.top() << '\n'; // 3
\end{lstlisting}

\subsection{pair 小根堆}

Dijkstra 常用：

\begin{lstlisting}
using P = pair<long long, int>; // distance, node
priority_queue<P, vector<P>, greater<P>> pq;
pq.push({0, 1});
\end{lstlisting}

pair 的默认比较先看 \Code{first}，再看 \Code{second}。

\section{常用算法函数}\label{sec:04-22}

\subsection{min、max、swap}

\begin{lstlisting}
int x = min(a, b);
int y = max(a, b);
swap(a, b);
\end{lstlisting}

\subsection{reverse}

\begin{lstlisting}
reverse(a.begin(), a.end());
reverse(s.begin(), s.end());
\end{lstlisting}

\subsection{unique 去重}

若要去掉所有重复值且允许改变原顺序，先排序。\Code{unique} 只消除相邻重复；要保留首次出现顺序，用集合配合一次扫描：

\begin{lstlisting}
sort(a.begin(), a.end());
a.erase(unique(a.begin(), a.end()), a.end());
\end{lstlisting}

\subsection{accumulate}

\begin{lstlisting}
long long sum = accumulate(a.begin(), a.end(), 0LL);
\end{lstlisting}

\Warn{第三个参数写 \Code{0LL}，否则可能按 \Code{int} 求和。}

\subsection{gcd}

\begin{lstlisting}
// Require a,b >= 0 and lcm(a,b) <= LLONG_MAX.
long long g = gcd(a, b);
long long l = (a == 0 || b == 0) ? 0 : a / g * b;
\end{lstlisting}

\Warn{最小公倍数先除后乘。}这样比 \Code{a * b / gcd(a,b)} 更不容易溢出。

\subsection{常用数学函数}

\begin{lstlisting}
double r = sqrt(x);      // square root
double y = pow(a, b);    // a^b, floating point
double f = floor(v);     // round down as floating point
double c = ceil(v);      // round up as floating point
long long z = llround(v); // nearest integer
\end{lstlisting}

提醒：

\begin{itemize}
  \item \Code{sqrt}、\Code{pow} 返回浮点数，不适合直接处理大整数精确幂。
  \item 整数平方尽量写 \Code{1LL * x * x}，不要写 \Code{pow(x, 2)}。
  \item 判断一个整数是否为完全平方数时，先取 \Code{sqrt}，再用整数乘法校验。
  \item 整数向上取整不要用 \Code{ceil(1.0*a/b)} 代替整数公式。
\end{itemize}

\begin{lstlisting}
bool isSquare(long long x) {
    if (x < 0) return false;
    long long r = sqrt((long double)x);
    while (r > 0 && r > x / r) r--;
    while ((r + 1) <= x / (r + 1)) r++;
    return r * r == x;
}
\end{lstlisting}

\section{lambda 入门}\label{sec:04-23}

lambda 是临时函数，排序时非常常用。

最基本写法：

\begin{lstlisting}
[](int a, int b) {
    return a < b;
}
\end{lstlisting}

用于排序：

\begin{lstlisting}
sort(a.begin(), a.end(), [](int x, int y) {
    return x > y;
});
\end{lstlisting}

使用外部变量：

\begin{lstlisting}
int mod = 10;
sort(a.begin(), a.end(), [mod](int x, int y) {
    return x % mod < y % mod;
});
\end{lstlisting}

引用捕获：

\begin{lstlisting}
int cnt = 0;
auto add = [&](int x) {
    cnt += x;
};
add(5);
\end{lstlisting}

考场建议：

\begin{itemize}
  \item 排序比较器中优先使用 \Code{const Type\&}，避免复制。
  \item 不熟悉捕获时，尽量不要写复杂 lambda。
  \item 复杂逻辑可以改成普通函数或先计算辅助字段。
\end{itemize}

\section{常见编译错误和原因}\label{sec:04-24}

\begin{longtable}{p{4cm}p{5cm}p{5cm}}
\toprule
报错或现象 & 常见原因 & 处理方法 \\
\midrule
\endfirsthead
\toprule
报错或现象 & 常见原因 & 处理方法 \\
\midrule
\endhead
\Code{expected ';'} & 少分号 & 看上一行结构体、变量定义、语句末尾。 \\
\Code{was not declared} & 变量未定义或作用域不对 & 检查拼写和定义位置。 \\
\Code{no matching function} & 函数参数类型不匹配 & 检查传参类型、引用、const。 \\
\Code{segmentation fault} & 越界、空容器取值、递归爆栈 & 查数组下标、\Code{front/top/back} 前是否非空。 \\
样例正确但大数据错 & 溢出或复杂度过高 & 检查 \Code{long long} 和复杂度。 \\
输出格式错误 & 多输出、少换行、空格不对 & 严格对照题目输出格式。 \\
\bottomrule
\end{longtable}

\section{C 和 C++ 混用提醒}\label{sec:04-25}

如果你使用 C++，建议全程使用 \Code{cin/cout} 或全程使用 \Code{scanf/printf}，不要混用。

如果确实使用 C 风格输入输出：

\begin{lstlisting}
int x;
scanf("%d", &x);

long long y;
scanf("%lld", &y);

printf("%d\n", x);
printf("%lld\n", y);
\end{lstlisting}

字符串：

\begin{lstlisting}
char s[1005];
scanf("%s", s);
\end{lstlisting}

初学时可以先熟悉 \Code{string}、\Code{vector} 和 \Code{map}，字符串和动态数组的存储会方便一些。

\section{考场 C++ 快速检查}\label{sec:04-26}

提交前快速看一遍：

\begin{itemize}
  \item 是否包含了头文件和 \Code{using namespace std;}。
  \item 是否开了快速输入输出。
  \item 是否混用了 \Code{cin/cout} 和 \Code{scanf/printf}。
  \item 求和、乘法、答案是否用了 \Code{long long}。
  \item vector 是否按正确大小初始化。
  \item 访问 \Code{front/top/back} 前是否判断非空。
  \item \Code{lower\_bound} 返回 \Code{end} 时是否处理。
  \item 排序比较器相等时是否返回 \Code{false} 或继续比较。
  \item 多组数据是否清空容器。
  \item 输出是否完全符合题意。
\end{itemize}
